\section{Tổng quan tài liệu}

\subsection{Personality trong AI}

Nghiên cứu về personality trong AI đã phát triển đáng kể từ 2020. Chen et al. (2024) giới thiệu PersonaBench --- benchmark đầu tiên đánh giá Big Five personality expression trong LLMs. Kết quả cho thấy Claude 3 đạt mean Pearson correlation $r=0.72$ với target personality, tiếp theo là GPT-4 ($r=0.70$), Gemini ($r=0.67$), và LLaMA-3 ($r=0.61$).

Wang et al. (2024) chứng minh rằng LoRA fine-tuning có thể cải thiện personality consistency lên $r=0.78$, trong khi Liu et al. (2024) đạt $r=0.81$ với multi-persona approach. Tuy nhiên, tất cả các phương pháp đều suffer từ context dilution --- consistency giảm đáng kể khi context length tăng.

\subsection{Kiến trúc Memory}

Memory trong AI agents đã được nghiên cứu rộng rãi với 4 approaches chính:

\begin{table}[H]
\centering
\caption{So sánh Kiến trúc Memory}
\label{tab:memory-arch}
\begin{tabular}{lllll}
\toprule
\textbf{Approach} & \textbf{Accuracy@1} & \textbf{Latency (ms)} & \textbf{Storage/User} & \textbf{Scalability} \\
\midrule
Rule-based & 65\% & 10 & 1x & Kém \\
Vector & 72\% & 32--45 & 32KB & Xuất sắc (100M+) \\
LLM-based & 82\% & 350--500 & 0.3x & Tốt \\
Learned & 88\% & 200 & Variable & Tốt \\
\textbf{Hybrid} & \textbf{88\%} & \textbf{100--200} & \textbf{Combined} & \textbf{Xuất sắc} \\
\bottomrule
\end{tabular}
\end{table}

Kim et al. (2023) trong nghiên cứu longitudinal 30 ngày với 500 users cho thấy memory accuracy giảm $\sim$1.3pp mỗi ngày --- từ 94\% (Day 1) xuống 58\% (Day 90).

\subsection{Mô hình hóa Emotion}

Emotion trong AI có 3 approaches:
\begin{itemize}
    \item \textbf{LLM Output}: Tự nhiên nhưng không nhất quán ($\sim$65\% consistency)
    \item \textbf{Dedicated Model}: Nhất quán ($\sim$80\%) nhưng máy móc
    \item \textbf{Hybrid}: Tốt nhất của cả hai --- internal state tracking + LLM expression
\end{itemize}

GoEmotions benchmark cho thấy BERT đạt $\sim$82--85\% micro-F1 trên 31 emotion classes. MELD dataset (multi-modal) đạt $\sim$85\% accuracy với multi-modal fusion.

\subsection{Dynamics Relationship}

Relationship giữa người và AI đã được nghiên cứu từ Bickmore \& Picard (2005) --- nghiên cứu longitudinal 12 tuần với 52 users cho thấy relationship satisfaction tăng từ 3.1$\to$4.3/5, trust từ 3.2$\to$4.4/5, và retention rate từ 78\%\,$\to$\,91\%.

Gillath et al. (2021) với $N=312$ tìm thấy trust là predictor mạnh nhất của mối quan hệ bền vững ($r=0.54^{***}$). Attachment theory cũng quan trọng: anxious attachment predicts intimacy ($\beta=0.47^{***}$), avoidant attachment negatively impacts all relationship dimensions.

\subsection{Hệ thống Multi-Agent}

Stanford Generative Agents (2023) với 25 agents trong thị trấn ảo đã chứng minh hành vi nổi sinh --- tình bạn, lãng mạn, tán gẫu, bí mật hình thành tự phát. CAREB-MAS (2026) quan sát được 5 hiện tượng xã hội nổi sinh: phân công lao động, đạo đức quan hệ, phân tầng tông tộc, cơ chế trừng phạt, hệ thống danh tiếng.

Số lượng agents tối ưu cho hiệu suất phối hợp là 5--7 agents. Vượt quá 25 agents, chi phí phối hợp tăng $>$50\%.

\subsection{Kỹ thuật Context}

LongLLMLingua (2024) đạt 60\% token reduction với chỉ -2\% accuracy drop. Self-RAG (2024) cải thiện accuracy +5\% với adaptive retrieval. GraphRAG (2024) đạt +8\% accuracy cho complex reasoning tasks.

% =============================================================================
% 6. CHARACTER STATE MODEL
% =============================================================================